\subsection{Unitary representations}

DEFINITION 3. --- Let G be a topological group and E a Hilbert space over K. A unitary representation of G in E is a continuous linear representation \(\varrho\) of G in E such that, for every g in G, the endomorphism \(\varrho(g)\) of E is a unitary endomorphism (EVT, V, p. 40) of E.

In other words, a unitary representation is an isometric representation in a Hilbert space. In particular, a unitary representation is bounded.

The trivial representation of a topological group in a Hilbert space is unitary. Every subrepresentation of a unitary representation is unitary. The restriction to a subgroup of a unitary representation is unitary.

DEFINITION 4. --- Let G be a topological group and let \(\varrho\) and \(\pi\) be unitary representations of G in Hilbert spaces E and F respectively. A G-invariant sesquilinear form on E × F is a continuous sesquilinear form q on E × F such that

\[
q (x, y) = q (\varrho (g) x, \pi (g) y)
\]

for all \(g \in \mathrm{G}\) and every \((x, y) \in \mathrm{E} \times \mathrm{F}\) .

The vector space of G-invariant sesquilinear forms on E × F is denoted \(\mathrm{Sesq}_{\mathrm{G}}(\varrho,\pi)\) or \(\mathrm{Sesq}_{\mathrm{G}}(\mathrm{E},\mathrm{F})\) .

Example. --- Let G be a topological group and let \(\varrho\) be a unitary representation of G in E. The scalar product \(q\) on E is a G-invariant sesquilinear form on \(\mathrm{E} \times \mathrm{E}\) .

Lemma 4. --- Let G be a topological group and let \(\varrho\) be a homomorphism of G into the unitary group U(E) of a Hilbert space E. Then \(\varrho\) is a unitary representation of G if and only if \(\varrho\) is continuous at the element e of G for the topology of simple convergence on U(E). It suffices that this property hold for every \(x\) in a total subset of E.

The condition is evidently necessary. It is sufficient by remark 2 of INT, VIII, p. 129, § 2, no. 1, since the image of \(\varrho\) is equicontinuous in \(\mathcal{L}(\mathrm{E})\) and since the continuity of the map \(g \mapsto \varrho(g)x\) at \(e\) implies its continuity on G.

Let G be a topological group. If \(\varrho_{1}\) and \(\varrho_{2}\) are unitary representations of G and if u belongs to \(\operatorname{Hom}_{\mathrm{G}}(\varrho_{1},\varrho_{2})\) , then \(u^{*}\) belongs to \(\operatorname{Hom}_{\mathrm{G}}(\varrho_{2},\varrho_{1})\) . Indeed, since \(\varrho_{1}\) and \(\varrho_{2}\) are unitary representations, one has for every \(g\in G\)

\[
\begin{array}{r} u ^ {*} \circ \varrho_ {2} (g) = u ^ {*} \circ \varrho_ {2} (g ^ {- 1}) ^ {*} = (\varrho_ {2} (g ^ {- 1}) \circ u) ^ {*} \\ = (u \circ \varrho_ {1} (g ^ {- 1})) ^ {*} = \varrho_ {1} (g) \circ u ^ {*}. \end{array}
\]

Lemma 5. --- Let G be a topological group and let \(\varrho\) be a unitary representation of G in a complex Hilbert space E. The space \(\mathrm{Hom}_{\mathrm{G}}(\varrho, \varrho)\) is a unital stellar subalgebra of \(\mathcal{L}(\mathrm{E})\) .

The preceding shows that \(\operatorname{Hom}_{\mathrm{G}}(\varrho,\varrho)\) is a unital and self-adjoint subalgebra of \(\mathcal{L}(\mathrm{E})\) . Since it is closed in \(\mathcal{L}(\mathrm{E})\) , it is a stellar subalgebra of \(\mathcal{L}(\mathrm{E})\) .

Lemma 6. --- Let \(\pi\) and \(\varrho\) be unitary representations of a topological group G in Hilbert spaces E and F respectively. Let D be a dense subspace of E which is stable under \(\pi\) . Let u be a closed partial operator from E to F with domain D such that \(u \circ \pi(g) = \varrho(g) \circ u\) for all \(g \in G\) . Then the domain of \(u^{*}\) is stable under \(\varrho\) and one has the relation \(u^{*} \circ \varrho(g) = \pi(g) \circ u^{*}\) for all \(g \in G\) .

Let \(g \in \mathrm{G}\) and \(x \in \mathrm{dom}(u^*)\) . For every \(y \in \mathrm{dom}(u)\) , we have

\[
\begin{array}{r l} & {\langle \varrho (g) x \mid u (y) \rangle = \langle x \mid \varrho (g ^ {- 1}) u (y) \rangle = \langle x \mid u (\pi (g ^ {- 1}) y) \rangle} \\ & {\qquad = \langle u ^ {*} (x) \mid \pi (g ^ {- 1}) y \rangle = \langle \pi (g) (u ^ {*} (x)) \mid y \rangle ,} \end{array}
\]

since \(\varrho(g)^{*} = \varrho(g)^{-1} = \varrho(g^{-1})\) . This proves that \(\varrho(g)x \in \mathrm{dom}(u^{*})\) and that \(u^{*}(\varrho(g)x) = \pi(g)(u^{*}(x))\) . In particular, the domain of \(u^{*} \circ \varrho(g)\) contains the domain of \(u^{*}\) , and \(\pi(g) \circ u^{*} \subset u^{*} \circ \varrho(g)\) .

But moreover, if \(x \in \text{dom}(u^{*} \circ \varrho(g))\) , then \(x = \varrho(g^{-1})(\varrho(g)x)\) belongs to \(\text{dom}(u^{*})\) by virtue of the preceding applied to \(g^{-1}\) . We conclude that \(u^{*} \circ \varrho(g) = \pi(g) \circ u^{*}\) .

PROPOSITION 1. --- Let \(\varrho_{1}\) and \(\varrho_{2}\) be unitary representations of a topological group G in Hilbert spaces E \(_{1}\) and E \(_{2}\) . The map from Hom \(_{G}(\varrho_{1},\varrho_{2})\) to Sesq \(_{G}(\varrho_{2},\varrho_{1})\) which associates to u the sesquilinear form \(q_{u}\) defined by \(q_{u}(x,y)=\langle x|u(y)\rangle\) is an isomorphism of vector spaces.

If \(u\) is a G-morphism, then one has

\[
\begin{array}{c} q _ {u} (\varrho_ {2} (g) x, \varrho_ {1} (g) y) = \langle \varrho_ {2} (g) x \mid u (\varrho_ {1} (g)y) \rangle \\ = \langle \varrho_ {2} (g) x \mid \varrho_ {2} (g) u (y) \rangle = \langle x \mid u (y) \rangle = q _ {u} (x, y) \end{array}
\]

for all \(g \in G\) and every \((x, y) \in \mathrm{E}_{2} \times \mathrm{E}_{1}\) , hence the indicated map is a linear map from \(\operatorname{Hom}_{\mathrm{G}}(\varrho_{1}, \varrho_{2})\) to \(\operatorname{Sesq}_{\mathrm{G}}(\varrho_{2}, \varrho_{1})\) . By EVT, V, p. 16, cor. 2, it is injective.

Conversely, let \(q\) be a G-invariant sesquilinear form and \(u\) the unique linear map from \(\mathrm{E}_1\) to \(\mathrm{E}_2\) such that \(q(x,y) = \langle x \mid u(y) \rangle\) for all \((x,y) \in \mathrm{E}_2 \times \mathrm{E}_1\) (loc. cit.). For every \(g\) in \(G\) and every \((x,y)\)

in \(\mathrm{E}_2\times \mathrm{E}_1\) , we have

\[
\begin{array}{c} \langle x | (\varrho_ {2} (g) \circ u \circ \varrho_ {1} (g ^ {- 1})) (y) \rangle = \langle \varrho_ {2} (g ^ {- 1}) x | u (\varrho_ {1} (g ^ {- 1}) y) \rangle \\ = q (\varrho_ {2} (g ^ {- 1}) x, \varrho_ {1} (g ^ {- 1}) y) = q (x, y) = \langle x | u (y) \rangle , \end{array}
\]

hence \(u = \varrho_{2}(g) \circ u \circ \varrho_{1}(g^{-1})\) . Consequently, \(u \in \operatorname{Hom}_{\mathrm{G}}(\varrho_{1}, \varrho_{2})\) . The proposition follows.

Let \(\varrho\) be a unitary representation of a topological group G in a Hilbert space E. If K = C, denote by \(\overline{\mathrm{E}}\) the conjugate space of E (EVT, V, p. 6). If K = R, set \(\overline{\mathrm{E}} = \mathrm{E}\) . The conjugate representation \(\overline{\varrho}\) is the representation of G in \(\overline{\mathrm{E}}\) defined by \(\overline{\varrho}(g) = \varrho(g)\) for all \(g \in \mathrm{G}\) . It is a unitary representation of G; by definition, a subspace of \(\overline{\mathrm{E}}\) is a subrepresentation of \(\overline{\mathrm{E}}\) if and only if it is a subrepresentation of E.

PROPOSITION 2. --- Let \(\varrho\) be a unitary representation of a topological group G in a Hilbert space E. Denote by u the isometric isomorphism of \(\overline{\mathrm{E}}\) into \(\mathrm{E}'\) which associates to \(x\) the linear form \(y\mapsto \langle x|y\rangle\) .

\begin{enumerate}
\def\labelenumi{\alph{enumi})}
\item
  The map u is an isomorphism of the conjugate representation \(\overline{\varrho}\) into the contragredient representation \(\breve{\varrho}\) ;
\item
  Endow E' with the Hilbert space structure obtained by transporting structure via u; the contragredient representation \(\breve{\varrho}\) is a unitary representation of G in E'.
\end{enumerate}

By EVT, V, p. 15, th. 3 and the following remark, the map \(u\) is an isometric isomorphism.

For every g in G, every x in \(\overline{E}\) and every y in E, we have

\[
\begin{array}{c} \langle y, (\breve {\varrho} (g) \circ u) (x) \rangle = \langle \varrho (g ^ {- 1}) y, u (x) \rangle \\ = \langle x | \varrho (g ^ {- 1}) y \rangle = \langle \varrho (g) x | y \rangle = \langle y, u (\overline {{\varrho}} (g) x) \rangle , \end{array}
\]

hence \(\breve{\varrho}(g) \circ u = u \circ \overline{\varrho}(g)\) , which proves \(a\) ; the assertion \(b\) ) follows immediately.

COROLLARY. --- Let \(\varrho\) be a unitary representation of a topological group G in a Hilbert space E. The following conditions are equivalent:

\begin{enumerate}
\def\labelenumi{(\roman{enumi})}
\item
  The representation \(\varrho\) is irreducible;
\item
  The contragredient representation \(\breve{\varrho}\) is irreducible;
\item
  The conjugate representation \(\overline{\varrho}\) is irreducible.
\end{enumerate}

This follows from Proposition 2, and from the remark preceding it concerning subrepresentations of \(\overline{E}\) .

PROPOSITION 3. --- Let \(\varrho\) be a unitary representation of a topological group G in a Hilbert space E. Let \(\pi\) be a subrepresentation of \(\varrho\), with space F. The orthogonal F° of F in E is a subrepresentation of E such that E = F ⊕ F° and F° is isomorphic to E/F.

The space F° is closed. For every \(x \in F^{\circ}\) , all \(g \in G\) and every \(y \in F\) , we have \(\langle \varrho(g)x \mid y \rangle = \langle x \mid \varrho(g^{-1})y \rangle = 0\) since \(\varrho\) is unitary and F is a subrepresentation. Therefore \(\varrho(g)x \in F^{\circ}\) .

The canonical projection of E onto E/F is a G-morphism, so its restriction to F° is a G-morphism from F° to E/F; by EVT, V, p. 13, it is an isometric isomorphism.

We shall also denote by \(\pi^{\circ}\) the representation of G in F°.

PROPOSITION 4. --- Let \(\varrho\) be a unitary representation of a topological group G in a Hilbert space E. A closed subspace F of E is a subrepresentation of \(\varrho\) if and only if the orthogonal projector p of E with image F is a G-morphism from \(\varrho\) to \(\varrho\) .

Suppose that F is a subrepresentation of \(\varrho\) . Let \(x \in E\) and let us write \(x = p(x) + y\) where \(y \in F^{\circ}\) . For every \(g\) in G, we have

\[
\varrho (g) x = \varrho (g) (p (x)) + \varrho (g) y,
\]

and since \(\varrho(g)(p(x))\) belongs to F and \(\varrho(g)y\) to F° (prop. 3), we have \(p(\varrho(g)x) = \varrho(g)(p(x))\) . Hence \(p\) belongs to \(\operatorname{Hom}_{\mathrm{G}}(\varrho, \varrho)\) .

Conversely, if \(p \in \operatorname{Hom}_{\mathrm{G}}(\varrho, \varrho)\) then \(1_{E} - p\) is a G-morphism, hence \(\mathrm{F} = \operatorname{Ker}(1_{\mathrm{E}} - p)\) is a subrepresentation of \(\varrho\) .

